% SPDX-License-Identifier: Apache-2.0
\section{A Tap, Sending Under a Predicate}
\label{sec:x-tap}

The stage of Section~\ref{sec:x-stage} waits: \code{until} a word
is offered and there is room, then receives and sends, and the
lowering makes the wait's condition the guard of every drive, the
\code{ready} of the receive and the \code{valid} of the send. A unit
that runs every cycle, as a processor core does, talks over its
channels differently. It receives with no wait before the receive,
which takes whatever is offered at every edge, so its \code{ready}
is its \code{valid}, the take itself, which is what the runtime
records as ready; and it sends under \code{when!}, so the channel's
\code{valid} is the arm's condition, or its negation in the
\code{else} arm, and the data is driven whatever the condition, since
data under a low \code{valid} means nothing. The tap takes every
word, counts them, and passes on the odd ones, counted too. This is
the form the Vreteno core's bus takes.

\begin{figure}[htbp]
\centering
\input{tap_timing.tex}
\caption{The tap: every word offered is taken, \code{inp\_ready}
with \code{inp\_valid}; the odd ones go out, \code{out\_valid}
the condition.}
\label{fig:x-tap}
\end{figure}

\lstinputlisting[caption={\code{ex\_tap.rs}. Run with \code{bazel run //lib/examples:ex\_tap}.},label={lst:x-tap}]{lib/examples/ex_tap.rs}

\lstinputlisting[style=ascii,caption={What \code{ex\_tap} printed when the build ran it: the words passed, then the lowering.},label={lst:x-tap-out}]{out_tap.txt}
